\documentclass[../../main.tex]{subfiles}
\begin{document}

{\bf [解]}
$e(\theta) = \cos\theta + i \sin\theta$ とし, $S$ の $n$ 個の要素を $a_k = e(\theta_k) \; (k=1, 2, \dots, n)$ とする. 
さらに (i) より $1$ が$S$ の要素だから $a_1 = 1$ とする. 
(ii)で $(z, w) = (1, 1)$ とすると $-1$ も $S$ の元であるから， $a_2 = -1$ とする. 以上最低でも二つの元があるから $n \geqq 2$ である.
よって題意の$n\le 4$に絞ると，$n=2, 3, 4$の3パターンについて調べれば良い．

\paragraph{1. $n=2$の時}
 $S = \{1, -1\}$ は (i), (ii) をみたす. よって十分である.

\paragraph{2. $n=3$の時}
$n=3$の時は$S=\pm$に加えてもう一つ$a_3$が要素として存在する．
(ii) で$(z, w) = (a_3, 1)$ とおくと,
 \begin{align}
    p = a_3 - 2 \cos\theta_3 = -\cos\theta_3 + i \sin\theta_3
 \end{align}
 が $S$ の元 $1, -1, a_3$ のいずれかに等しい.
 \begin{align}
 \begin{cases}
  p = 1 \iff \sin\theta_3 = 0 \land -\cos\theta_3 = 1 \iff a_3 = -1 \\
  p = -1 \iff a_3 = 1 \\
  p = a_3 \iff a_3 = \pm i
 \end{cases}
\end{align}
 だから, $a_3 \neq 1, -1$ とあわせて, $a_3 = \pm i$ が必要だが, $S = \{1, -1, i\}, \{1, -1, -i\}$ はいずれも (ii)で $(z, w) = (i, i)$ or $(-i, -i)$ とすればそれぞれ$p=-i, i$も$S$の元となってしまい不適である．

\paragraph{3. $n=4$の時}

 $2.$から $a_3 = \pm i$ が必要である. 
まず $a_3 = i$ とすると (ii)で $(z, w) = (i, i)$ として, $-i$ も $S$ の元である. よって $S = \{\pm 1, \pm i\}$ は (i), (ii) をみたし十分である．
$a_3 = -i$ の時も同じく, $(z, w) = (-i, -i)$ として $S = \{\pm 1, \pm i\}$ を得る.

以上3つの場合分けから求める$S$は 
\begin{align}
 S = \{\pm 1\}, \{\pm 1, \pm i\} 
\end{align}
である．


\end{document}
